\section{Digital-development results and experimental assessment}
\label{sec:results}
\label{sec:deposited_assessment}
\subsection{Two complementary development tasks}
The low-speed case begins with a predefined structure and discovers context-dependent commands. The high-speed case combines structure and dependency discovery. Proprietary structure identifiers, slot assignments and command sequences are withheld. Both pass through coupled fitting, integer-command realization, deployment and deposited-position assessment. A best-found structure is not claimed to be globally optimal.

\input{sections/framework_evidence}

\subsection{Operating points and paired measurements}
Computational design was followed by deployment of the resulting DAC table on the printer, physical printing and independent deposited-coordinate extraction. The low-speed case illustrates dependency discovery on a predefined structure. The high-speed case illustrates joint structure and dependency discovery. The main text describes the cases by operating point; the corresponding reproduction inputs are provided with the implementation.

According to the clarified acquisition provenance, ten repeated prints of the same complete uppercase A--Z message were analyzed from images for each raster. Corresponding dot coordinates were averaged across the ten measurements, and the supplied arrays contain these mean positions. There are ten repeated message prints per case; the 26 letters are not treated as independent replicate prints. The collected samples cover uppercase A--Z at 75 kHz and 60 requested horizontal DPI. Substrate speed is 50 m/min for the low-speed case and 100 m/min for the high-speed case. Each supplied coordinate dataset contains 1,725 mean dot positions. Coordinates are in millimeters. A nominal deposited diameter of $528\um$ is used only for footprint visualization; localization error is computed from dot centers.

The low-speed reference is a fresh native-75-kHz physical replay of the unchanged exported integer-DAC table. The high-speed reference uses its native-75-kHz repaired integer-DAC replay. The amplifier and induction sequence are evaluated at that frequency before integrating particle flight. No timing projection substitutes for physical integration, and no commands are refitted to the measured positions. The cached message reconstruction uses previous/current-stroke references with blank older history and no next-stroke emission; it is not an uninterrupted whole-message simulation. The original campaigns also used different native droplet-property records, so this is not a controlled causal comparison of structure alone.

Dot identities are matched by message line, column and logical landing label, rather than nearest-neighbor reassignment. All measured dots are retained. Input hashes, frozen simulation references and derived metrics are preserved in the accompanying repository. The comparison uses a versioned experimental coordinate dataset and fixed simulation predictions; superseded coordinate versions are retained separately and are not pooled as independent acquisitions.

\subsection{Character reconstruction from deposited coordinates}
Figures~\ref{fig:coordinate_low}--\ref{fig:coordinate_high} reconstruct the uppercase glyph samples directly from the paired simulation and measured center coordinates. They are coordinate-based reconstructions, not photographs or newly generated font targets. All 1,725 matched dots per case are retained. Each measured/simulated pair uses identical millimeter scales, bounds and a nominal $528\um$ circular footprint for visual context; actual measured dot diameters are unavailable. The experimental reconstruction depicts mean dot positions across ten prints, rather than any one physical print. Overlay panels show the matched centers at their original relative positions, with no translation, rotation, scale, shear or per-character registration fitted. Letters are displayed in alphabet order; the plotted experimental centers represent the ten-print mean message.

\begin{figure}[p]\centering\includegraphics[width=\linewidth,height=0.78\textheight,keepaspectratio]{figures/coordinate_print_low.pdf}\caption{Low-speed case: simulation and measured-print reconstruction from deposited coordinates, followed by matched-center overlays. The same coordinates used for the quantitative errors are shown without an alignment fit. Nominal circle sizes illustrate footprint only.}\label{fig:coordinate_low}\end{figure}
\begin{figure}[p]\centering\includegraphics[width=\linewidth,height=0.78\textheight,keepaspectratio]{figures/coordinate_print_high.pdf}\caption{High-speed case: simulation, measured-print reconstruction and matched-center overlays, with equal millimeter axes and all uppercase glyph samples retained.}\label{fig:coordinate_high}\end{figure}
\begin{figure}[htbp]\centering\includegraphics[width=\linewidth]{figures/coordinate_print_closeups.pdf}\caption{Selected character-center overlays at unit displacement scale. Blue dots are simulated centers, orange crosses are measured centers, and gray segments connect matched identities. One identical plotting-origin subtraction is applied to both coordinate sets within each glyph; it does not alter their residuals. These shape examples supplement the complete A--Z views and all-dot metrics.}\label{fig:coordinate_closeups}\end{figure}

\subsection{Prediction accuracy and target quality}
Let $\vect p_i=(y_i,x_i)$ denote the simulated substrate-direction and deflection coordinates, and let $\vect m_i$ denote the corresponding measured position. Prediction residuals and pooled root-mean-square errors are
\begin{equation}
 \vect e_i=\vect m_i-\vect p_i,\quad
 E_a=\sqrt{\frac1n\sum_{i=1}^n e_{i,a}^2},\quad
 E_{2D}=\sqrt{\frac1n\sum_{i=1}^n\lVert\vect e_i\rVert_2^2}.
\end{equation}
These are errors against physical deposited coordinates, distinct from an optimized simulation's residual against an ideal geometric target. A nearly ideal computational print is not itself physical validation.

Table~\ref{tab:deposited_errors} reports raw-coordinate comparisons without fitting translation, rotation, scale, shear, or per-character corrections. The simulator uses one phase defined from its own blank/full-stroke reference, not fitted to measured dots. Whole-file translation and rigid registrations are available separately in the internal sensitivity analysis; those are in-sample nuisance-parameter fits, not additional independent experiments.

\begin{table}[htbp]\centering\small
\caption{Measured-minus-simulated deposited-coordinate errors at the two tested operating points. All errors are in micrometers; P95 refers to radial residual magnitude.}\label{tab:deposited_errors}
\begin{tabular}{lrrrrr}\toprule
Case & Dots & Horizontal RMS & Deflection RMS & 2D RMS & 2D P95\\\midrule
@@DEPOSITED_ERROR_ROWS@@
\bottomrule\end{tabular}
\end{table}

\begin{figure}[htbp]\centering
\includegraphics[width=\linewidth]{figures/experimental_residual_distributions.pdf}
\caption{Unregistered residual distributions relative to the experimental mean positions. The cases represent different design tasks and operating points.}
\end{figure}

\begin{figure}[htbp]\centering\includegraphics[width=\linewidth]{figures/framework_glyph_assessment.pdf}\caption{Per-character deflection RMSE against the experimental ten-print mean. Commands and coordinate registration are fixed across the assessment.}\label{fig:per_character}\end{figure}

@@EXPERIMENTAL_INTERPRETATION@@

\subsection{Experimental scope and confidentiality}
The released assessment contains 1,725 mean dot positions per raster, obtained from ten repeated A--Z prints at each operating point. Additional laboratory records, including raw images and individual replicate measurements, remain confidential and are outside the public evidence package. The reported RMSE therefore measures disagreement with the released mean coordinates; it does not estimate single-print variability or a repeatability confidence interval.

This assessment tests deposited-position prediction and the computational-to-printer workflow for the two reported configurations. It does not independently identify each force closure or assess intermediate trajectories, contact events or collection margins. Noncoordinate fields retained in the export schema, such as flight time and gutter diagnostics, are simulation metadata rather than additional measured labels. The experimental conclusions are restricted to the released coordinates and operating points.

@@PHYSICS_ABLATION@@
